\documentclass[11pt]{article}

\usepackage[margin=1in]{geometry}
\usepackage{amsmath,amssymb,amsthm}
\usepackage{hyperref}
\textheight9in
\textwidth6.75in

\newtheorem{fact}{Fact}

% Notation shortcuts
\newcommand{\R}{\mathbb{R}}
\newcommand{\E}{\mathbb{E}}
\newcommand{\PP}{\mathbb{P}}
\newcommand{\Var}{\operatorname{Var}}
\newcommand{\opnorm}[1]{\left\lVert #1 \right\rVert_{\mathrm{op}}}
\newcommand{\norm}[1]{\left\lVert #1 \right\rVert}
\newcommand{\inner}[2]{\langle #1, #2 \rangle}
\newcommand{\lam}{\lambda}

\title{Assumed Background Facts for the Lean Formalization\\
of ``Estimation Error in Latent High-Dimensional Factor Models''}
\author{}
\date{}

\begin{document}
\maketitle

\noindent
The three statements below are the only external inputs the formalization
must assume. Each is (i)~classical, appearing in standard texts, (ii)~absent
from the current Mathlib library, and (iii)~\emph{not} proved in the paper. The
paper quotes each by name, and every other result in the paper
is proved there and is to be derived in Lean from these three facts together with material already
in Mathlib.

Throughout, $(\Omega,\mathcal F,\PP)$ is a probability space; $A,A',A^{(p)}$
denote real symmetric $m\times m$ matrices with eigenvalues listed in weakly
decreasing order $\lam_1\ge\cdots\ge\lam_m$; $\opnorm{\cdot}$ is the operator
(spectral) norm.

\bigskip

\begin{fact}[Kolmogorov strong law for independent, non-identically-distributed summands]
\label{Fact1}
Let $\{X_i\}_{i\ge 1}$ be independent real random variables on
$(\Omega,\mathcal F,\PP)$ with finite variances satisfying Kolmogorov's series
condition
\[
  \sum_{i=1}^{\infty}\frac{\Var(X_i)}{i^{2}} < \infty .
\]
Then we have
\[
  \frac{1}{p}\sum_{i=1}^{p}\bigl(X_i-\E[X_i]\bigr)
  \;\xrightarrow[p\to\infty]{}\; 0
  \qquad \PP\text{-almost surely.}
\]
\end{fact}


\medskip
\noindent\textbf{Reference.} V.~V.~Petrov, \emph{Sums of Independent Random
Variables}, translated by A.~A.~Brown, Ergebnisse der Mathematik und ihrer
Grenzgebiete, Band 82, Springer-Verlag, 1975 (the strong law of large numbers
for sums of independent, not necessarily identically distributed, random
variables).

\bigskip

\begin{fact}[Weyl's eigenvalue perturbation inequality]
\label{Fact2}
For real symmetric matrices $A, A' \in \R^{m\times m}$,
\[
  \bigl\lvert \lam_i(A') - \lam_i(A) \bigr\rvert \;\le\; \opnorm{A' - A},
  \qquad i = 1,\dots,m .
\]
Consequently, if $A^{(p)}\to A$ (equivalently, since $m$ is fixed, entrywise),
then $\lam_i\bigl(A^{(p)}\bigr)\to\lam_i(A)$ for every $i=1,\dots,m$.
\end{fact}



\medskip
\noindent\textbf{Reference.} R.~Bhatia, \emph{Matrix Analysis}, Graduate Texts
in Mathematics 169, Springer, 1997, Cor.~III.2.6.

\bigskip

\begin{fact}[Eigenvector continuity for a simple eigenvalue]
\label{Fact3}
Let $A^{(p)}, A \in \R^{m\times m}$ be real symmetric with $A^{(p)}\to A$.
Suppose $\mu$ is a \emph{simple} eigenvalue of $A$ with unit eigenvector $v$
(every unit vector $w$ with $Aw=\mu w$ satisfies $w=\pm v$). Then there exist
$\mu^{(p)}\in\R$ and unit vectors $v^{(p)}\in\R^{m}$ such that, for all
sufficiently large $p$,
\[
  A^{(p)} v^{(p)} = \mu^{(p)} v^{(p)},
  \qquad \norm{v^{(p)}}_2 = 1,
  \qquad \inner{v^{(p)}}{v} \ge 0,
\]
and $\mu^{(p)}\to\mu$, $v^{(p)}\to v$ as $p\to\infty$.
\end{fact}


\medskip
\noindent\textbf{Reference.} C.~Davis and W.~M.~Kahan, ``The rotation of
eigenvectors by a perturbation.~III,'' \emph{SIAM Journal on Numerical
Analysis} \textbf{7} (1970), no.~1, 1--46.

\end{document}
